\begin{abstract}
Agentic skills combine natural language instructions with scripts, tools,
files, and external services, allowing security relevant behavior to span
multiple artifacts and stages of execution. Existing scanners can identify
suspicious capabilities, but capability alone does not establish a realized
security violation. A protected file and a network endpoint may coexist
without data flow, and an observed flow may still be authorized. Security
analysis must establish what behavior is expressed by the skill, what
source-to-sink relation is realized during execution, and what security meaning
that relation has under policy.

\system addresses this problem through evidence constrained security
adjudication. Static analysis constructs grounded instruction provenance over
security relevant sources, actions, destinations, and supporting artifact
spans. Bounded execution follows protected sources through LLM context, tool
interactions, files, processes, and network carriers. Cross layer alignment
connects compatible instruction and runtime relations while preserving their
evidence origin. A closure contract determines whether the recovered evidence
supports a security relevant path, and policy evaluation classifies supported
behavior as violating, allowed, or unresolved. Coverage records where runtime
evidence ends, keeping incomplete executions explicit in the resulting
assessment.

We evaluate \system on ProvBench, a 776-case benchmark covering confirmed
violations, benign lookalikes, trusted allowed flows, incomplete executions,
and 142 counterfactual pairs. Static screening reaches 1.000 recall and 0.917
F1, while full adjudication achieves 0.720 F1 for violation closure. Evidence
closure recovers supported source-to-sink relations for 61.6\% of confirmed
violations with no false evidence closures in the evaluated corpus, and the
full pipeline produces no false positives on the 179 benign lookalikes. We
further audit 6,000 public skills, applying bounded dynamic replay to 47.8\% of
the corpus and using recovered provenance to guide manual verification. These
results show that instruction and runtime provenance can move agentic skill
security from suspicious capability screening toward reviewable,
evidence-supported security decisions.
\end{abstract}